\begin{pastexam}[label={px:ncku-113-prog-10}]{成大資訊聯招 113 學年度 第 10 題}
（單選）

\textbf{[5\%]} A binary trie can be transformed to Patricia. Considering the following binary trie,
which of the Patricia is a correct transformation result?
\begin{center}
\begin{tikzpicture}[x=0.0145cm,y=-0.0145cm,b/.style={ellipse,draw,fill=white,minimum width=1.05cm,minimum height=0.56cm,inner sep=0pt},e/.style={rectangle,draw,fill=white,font=\sffamily\small,inner sep=2.5pt}]
\node[b] (r) at (670,1012) {};
\node[b] (l1) at (587,1067) {};
\node[b] (r1) at (754,1067) {};
\node[b] (l2) at (532,1124) {};
\node[b] (r2) at (829,1124) {};
\node[b] (l3) at (609,1182) {};
\node[b] (r3) at (786,1183) {};
\node[e] (k1000) at (706,1122) {1000};
\node[e] (k0000) at (475,1188) {0000};
\node[e] (k0010) at (563,1241) {0010};
\node[e] (k0011) at (643,1241) {0011};
\node[e] (k1100) at (741,1241) {1100};
\node[e] (k1101) at (829,1241) {1101};
\draw (r)--(l1);
\draw (r)--(r1);
\draw (l1)--(l2);
\draw (l2)--(k0000);
\draw (l2)--(l3);
\draw (l3)--(k0010);
\draw (l3)--(k0011);
\draw (r1)--(k1000);
\draw (r1)--(r2);
\draw (r2)--(r3);
\draw (r3)--(k1100);
\draw (r3)--(k1101);
\end{tikzpicture}
\end{center}

\begin{enumerate}[label=(\Alph*)]
\item \begin{tikzpicture}[x=0.0075cm,y=-0.0075cm,baseline=(current bounding box.north),p/.style={ellipse,draw,fill=white,minimum width=1.3cm,minimum height=0.64cm,inner sep=0pt,font=\sffamily\small},k/.style={font=\sffamily\bfseries\large,inner sep=0pt},a/.style={->,line width=0.9pt,>={Straight Barb[length=3.5pt,width=5pt]}}]
\node[p] (p0000) at (632,133) {0000};
\node[p] (p1000) at (385,218) {1000};
\node[p] (p0011) at (200,347) {0011};
\node[p] (p1100) at (548,340) {1100};
\node[p] (p0010) at (347,480) {0010};
\node[p] (p1101) at (684,478) {1101};
\node[k] at (665,65) {0};
\node[k] at (372,150) {1};
\node[k] at (170,275) {3};
\node[k] at (572,268) {2};
\node[k] at (375,412) {4};
\node[k] at (705,408) {4};
\draw[a] (p0000)--(p1000);
\draw[a] (p1000)--(p0011);
\draw[a] (p1000)--(p1100);
\draw[a] (p0011)--(p0010);
\draw[a] (p1100)--(p1101);
\draw[a] (112,347) .. controls (50,250) and (50,110) .. (160,75) .. controls (300,35) and (500,45) .. (628,88);
\draw[a] (462,340) .. controls (415,330) and (392,300) .. (387,263);
\draw[a] (262,492) .. controls (205,580) and (205,385) .. (290,447);
\draw[a] (408,512) .. controls (420,590) and (485,570) .. (472,470) .. controls (462,400) and (380,355) .. (290,347);
\draw[a] (615,508) .. controls (590,580) and (515,585) .. (500,500) .. controls (492,450) and (488,410) .. (490,372);
\draw[a] (762,495) .. controls (815,580) and (815,385) .. (745,443);
\end{tikzpicture}
\item \begin{tikzpicture}[x=0.0075cm,y=-0.0075cm,baseline=(current bounding box.north),p/.style={ellipse,draw,fill=white,minimum width=1.3cm,minimum height=0.64cm,inner sep=0pt,font=\sffamily\small},k/.style={font=\sffamily\bfseries\large,inner sep=0pt},a/.style={->,line width=0.9pt,>={Straight Barb[length=3.5pt,width=5pt]}}]
\node[p] (p0010) at (662,140) {0010};
\node[p] (p1000) at (412,225) {1000};
\node[p] (p0000) at (228,355) {0000};
\node[p] (p1100) at (580,348) {1100};
\node[p] (p0011) at (375,490) {0011};
\node[p] (p1101) at (718,485) {1101};
\node[k] at (695,65) {0};
\node[k] at (400,158) {1};
\node[k] at (195,288) {2};
\node[k] at (605,275) {2};
\node[k] at (405,425) {3};
\node[k] at (738,415) {3};
\draw[a] (p0010)--(p1000);
\draw[a] (p1000)--(p0000);
\draw[a] (p1000)--(p1100);
\draw[a] (p0000)--(p0011);
\draw[a] (p1100)--(p1101);
\draw[a] (303,512) .. controls (180,600) and (80,520) .. (70,330) .. controls (62,120) and (300,45) .. (598,100);
\draw[a] (494,348) .. controls (445,338) and (420,308) .. (415,270);
\draw[a] (143,368) .. controls (95,470) and (95,255) .. (165,325);
\draw[a] (440,518) .. controls (465,600) and (545,545) .. (464,488);
\draw[a] (650,514) .. controls (610,600) and (540,570) .. (520,382);
\draw[a] (795,505) .. controls (850,590) and (850,385) .. (780,448);
\end{tikzpicture}
\item \begin{tikzpicture}[x=0.0075cm,y=-0.0075cm,baseline=(current bounding box.north),p/.style={ellipse,draw,fill=white,minimum width=1.3cm,minimum height=0.64cm,inner sep=0pt,font=\sffamily\small},k/.style={font=\sffamily\bfseries\large,inner sep=0pt},a/.style={->,line width=0.9pt,>={Straight Barb[length=3.5pt,width=5pt]}}]
\node[p] (p0000) at (768,137) {0000};
\node[p] (p1000) at (505,230) {1000};
\node[p] (p0011) at (310,372) {0011};
\node[p] (p1100) at (680,360) {1100};
\node[p] (p0010) at (222,515) {0010};
\node[p] (p1101) at (630,522) {1101};
\node[k] at (802,65) {0};
\node[k] at (493,160) {1};
\node[k] at (275,297) {3};
\node[k] at (707,285) {2};
\node[k] at (165,443) {4};
\node[k] at (570,440) {4};
\draw[a] (p0000)--(p1000);
\draw[a] (p1000)--(p0011);
\draw[a] (p1000)--(p1100);
\draw[a] (p0011)--(p0010);
\draw[a] (p1100)--(p1101);
\draw[a] (222,372) .. controls (150,250) and (150,60) .. (450,55) .. controls (600,50) and (700,65) .. (762,92);
\draw[a] (594,360) .. controls (545,350) and (515,318) .. (508,275);
\draw[a] (137,530) .. controls (55,640) and (55,395) .. (158,483);
\draw[a] (290,545) .. controls (340,640) and (450,600) .. (450,480) .. controls (450,420) and (435,385) .. (400,372);
\draw[a] (548,540) .. controls (490,640) and (490,410) .. (566,490);
\draw[a] (695,552) .. controls (740,625) and (820,585) .. (820,470) .. controls (820,410) and (805,372) .. (770,360);
\end{tikzpicture}
\item \begin{tikzpicture}[x=0.0075cm,y=-0.0075cm,baseline=(current bounding box.north),p/.style={ellipse,draw,fill=white,minimum width=1.3cm,minimum height=0.64cm,inner sep=0pt,font=\sffamily\small},k/.style={font=\sffamily\bfseries\large,inner sep=0pt},a/.style={->,line width=0.9pt,>={Straight Barb[length=3.5pt,width=5pt]}}]
\node[p] (p0000) at (768,137) {0000};
\node[p] (p1000) at (505,230) {1000};
\node[p] (p0011) at (310,372) {0011};
\node[p] (p1100) at (680,360) {1100};
\node[p] (p0010) at (222,515) {0010};
\node[p] (p1101) at (630,522) {1101};
\node[k] at (802,65) {0};
\node[k] at (493,160) {1};
\node[k] at (275,297) {3};
\node[k] at (707,285) {2};
\node[k] at (165,443) {3};
\node[k] at (570,440) {3};
\draw[a] (p0000)--(p1000);
\draw[a] (p1000)--(p0011);
\draw[a] (p1000)--(p1100);
\draw[a] (p0011)--(p0010);
\draw[a] (p1100)--(p1101);
\draw[a] (222,372) .. controls (150,250) and (150,60) .. (450,55) .. controls (600,50) and (700,65) .. (762,92);
\draw[a] (594,360) .. controls (545,350) and (515,318) .. (508,275);
\draw[a] (137,530) .. controls (55,640) and (55,395) .. (158,483);
\draw[a] (290,545) .. controls (340,640) and (450,600) .. (450,480) .. controls (450,420) and (435,385) .. (400,372);
\draw[a] (548,540) .. controls (490,640) and (490,410) .. (566,490);
\draw[a] (695,552) .. controls (740,625) and (820,585) .. (820,470) .. controls (820,410) and (805,372) .. (770,360);
\end{tikzpicture}
\item None of the above
\end{enumerate}
\end{pastexam}

\begin{pastsoln}
答案是 (A)。

位元從左邊數起，第 $1$ 到第 $4$ 位。 Patricia 的定義見資結書第~11.5.3~節：先把二元 trie 壓縮
（拿掉只有一個子節點的分支節點），每個分支節點記下它看第幾個位元（bitNumber），再把鍵放進分支節點，
加一個 bitNumber 為 $0$ 的標頭節點；指向元素的指標改成指回放著那個鍵的節點（往回的指標）。

\textbf{1. 壓縮二元 trie}。 題目的 trie 裡，左邊第 $2$ 層（所有鍵第 $2$ 位都是 $0$）和右邊第 $3$ 層的
分支節點只有一個子節點，拿掉。 每個分支節點的 bitNumber 是它子樹裡的鍵第一個不全相同的位元：
\begin{itemize}
\item 根：六個鍵第 $1$ 位就不同，bit $1$。
\item 左子樹 $\{0000, 0010, 0011\}$：第 $2$ 位都是 $0$，第 $3$ 位才不同，bit $3$；左邊是 $0000$，右邊
$\{0010, 0011\}$ 在第 $4$ 位分開，bit $4$。
\item 右子樹 $\{1000, 1100, 1101\}$：第 $2$ 位不同，bit $2$；左邊是 $1000$，右邊 $\{1100, 1101\}$ 在
第 $4$ 位分開，bit $4$。
\end{itemize}
所以 Patricia 除了標頭（bit $0$）之外，bitNumber 由上往下是 $1$；$3$、$2$；$4$、$4$。 (B) 有兩個 bit
$2$、兩個 bit $3$，(D) 的兩個最下層節點是 bit $3$，都不對。 剩下 (A)、(C)，形狀一樣，差在左右。

\textbf{2. 左右要對}。 bit $3$ 節點的右邊（第 $3$ 位是 $1$）才是 $\{0010, 0011\}$；bit $2$ 節點的右邊
（第 $2$ 位是 $1$）才是 $\{1100, 1101\}$。 所以 bit $4$ 的兩個節點都應該是\textbf{右}子節點。 (C)
把它們畫成左子節點；而且 (C) 裡 $1101$ 指向自己的是左指標，但 $1101$ 的第 $4$ 位是 $1$，搜尋
$1101$ 會走右指標，到不了自己。 (C) 錯。

\textbf{3. 檢查 (A)}。 用資結書虛擬碼~11.11 的搜尋，對每個鍵從標頭的左指標開始，照 bitNumber 指定的位元
往下走，直到走了一個往回的指標（指到 bitNumber 不比目前大的節點）：
\begin{center}
\small
\begin{tabular}{lll}
\toprule
鍵 & 路徑（括號是看的位元和它的值） & 停在 \\
\midrule
0000 & 1000（第 1 位 0）$\to$ 0011（第 3 位 0）$\to$ 往回 & 0000 \\
0010 & 1000（第 1 位 0）$\to$ 0011（第 3 位 1）$\to$ 0010（第 4 位 0）$\to$ 往回 & 0010 \\
0011 & 1000（第 1 位 0）$\to$ 0011（第 3 位 1）$\to$ 0010（第 4 位 1）$\to$ 往回 & 0011 \\
1000 & 1000（第 1 位 1）$\to$ 1100（第 2 位 0）$\to$ 往回 & 1000 \\
1100 & 1000（第 1 位 1）$\to$ 1100（第 2 位 1）$\to$ 1101（第 4 位 0）$\to$ 往回 & 1100 \\
1101 & 1000（第 1 位 1）$\to$ 1100（第 2 位 1）$\to$ 1101（第 4 位 1）$\to$ 往回 & 1101 \\
\bottomrule
\end{tabular}
\end{center}
六個鍵都停在放著自己的節點，(A) 的每一個往回的指標都正確：$0011$ 的左指標指回標頭 $0000$，$0010$
的左指標指自己、右指標指 $0011$，$1100$ 的左指標指 $1000$，$1101$ 的左指標指 $1100$、右指標指自己。
所以 (A) 是正確的 Patricia。 （哪個鍵放在哪個節點和插入順序有關，例如依 $0000, 1000, 0011, 1100,
0010, 1101$ 的順序插入就會得到 (A)；不同順序可能得到別的放法，但選項裡只有 (A) 的指標全部正確。）
\end{pastsoln}
